% monopol_spektral_analyse.tex
% Standalone-Abschnitt für das EABC-Monopol-Preprint (arXiv).
% Einbindung im Hauptdokument: % monopol_spektral_analyse.tex
% Standalone-Abschnitt für das EABC-Monopol-Preprint (arXiv).
% Einbindung im Hauptdokument: % monopol_spektral_analyse.tex
% Standalone-Abschnitt für das EABC-Monopol-Preprint (arXiv).
% Einbindung im Hauptdokument: % monopol_spektral_analyse.tex
% Standalone-Abschnitt für das EABC-Monopol-Preprint (arXiv).
% Einbindung im Hauptdokument: \input{monopol_spektral_analyse}
% Voraussetzungen: amsmath, booktabs (ggf. hyperref für \ref)

\input{monopol_preprint_results.txt}

% Modellkonstanten (nicht laufzeitabhängig)
\newcommand{\monopolQtot}{1796}
\newcommand{\monopolFactor}{1036.92} % 1796/\sqrt{3}
\newcommand{\monopolPressureNorm}{0.032292} % \monopolPressure / \monopolN

% Konvergenzdaten (Teilstichproben, gleicher Lauf wie oben)
\newcommand{\convFiftyVar}{0.172806}
\newcommand{\convFiftyDelta}{0.068616}
\newcommand{\convFiftyMean}{0.992296}
\newcommand{\convHundredVar}{0.155540}
\newcommand{\convHundredDelta}{0.051350}
\newcommand{\convHundredMean}{0.995738}
\newcommand{\convTwoHundredVar}{0.160372}
\newcommand{\convTwoHundredDelta}{0.056182}
\newcommand{\convTwoHundredMean}{0.998893}

% Diagnosewerte der fehlerhaften Hauptterm-Pipeline (Referenzlauf, $t\in[5000,5140]$)
\newcommand{\badMeanFixed}{2.9586}
\newcommand{\badVarFixed}{7.33}
\newcommand{\badMeanAdapt}{2.9962}
\newcommand{\badVarAdapt}{7.88}
\newcommand{\badCountInterval}{52}
\newcommand{\goodCountInterval}{150}

%=============================================================================
\section{Methodische Richtigstellung zur numerischen Nullstellen-Asymptotik bei hohen Parametern}
\label{sec:methodische-richtigstellung}
%=============================================================================

In der numerischen Verifikation des EABC-Monopolmodells im hochenergetischen Bereich
($t \approx 5000$) zeigt sich eine methodische Divergenz, die für jede spektrale
Auswertung zuerst geklärt werden muss.
Eine Auswertung der Riemann-Siegel-Z-Funktion, die nur den asymptotischen Hauptterm
verwendet und den Restterm $R(t)$ vernachlässigt, führt bei der Nullstellensuche
über Vorzeichenwechsel zu systematischen Lücken in der Stichprobe.

\begin{table}[htbp]
\centering
\caption{Vergleich numerischer Approximationsmethoden im Intervall $t \in [5000,\,5140]$.
Die ersten beiden Zeilen verwenden den unvollständigen Hauptterm; die dritte Zeile die
vollständige Implementierung \texttt{mpmath.siegelz} mit adaptiver Schrittweite
und \texttt{brentq}-Verfeinerung (Stichprobe $N=50$).}
\label{tab:siegel_vergleich}
\begin{tabular}{lccc}
\toprule
\textbf{Methode} &
\textbf{Mittlerer norm.\ Abstand} &
\textbf{Varianz $\sigma^2$} &
\textbf{Gefundene Nullstellen} \\
\midrule
Hauptterm, Schrittweite $0{,}5$ &
$\badMeanFixed$ & $\badVarFixed$ & $\sim\badCountInterval$ \\
Hauptterm, adaptive Schrittweite &
$\badMeanAdapt$ & $\badVarAdapt$ & $\sim\badCountInterval$ \\
\texttt{mpmath.siegelz}, adaptiv &
$\convFiftyMean$ & $\convFiftyVar$ & $\sim\goodCountInterval$ \\
\bottomrule
\end{tabular}
\end{table}

Tabelle~\ref{tab:siegel_vergleich} zeigt: Eine rein adaptive Schrittweite behebt das
Problem nicht, solange $R(t)$ fehlt---es werden dieselben $\sim\badCountInterval$ Nullstellen
gefunden wie mit fester Schrittweite $0{,}5$.
Der scheinbare Faktor $\sim 3$ im mittleren normierten Abstand ist damit ein
Aliasing-Effekt der Unterabtastung, kein spektraler Beitrag des Monopols.

Erst unter Einbeziehung der vollständigen Riemann-Siegel-Asymptotik stellt sich
Weyl-konformes Verhalten mit mittlerem Abstand $\approx 1$ ein.
Eine Modifikation der globalen Zählfunktion der Form
$N_{\mathrm{EABC}}(T) \approx \tfrac{1}{3}\,N_{\mathrm{standard}}(T)$
--- oder jede daraus abgeleitete ``Drittel-Symmetrie'' der Nullstellendichte ---
ist ohne unabhängige Evidenz nicht haltbar und wird verworfen.

%=============================================================================
\section{Spektrale Abweichungsanalyse und Monopol-Kopplung}
\label{sec:spektrale-abweichung}
%=============================================================================

\subsection{Statistisches Setup}

Nach der methodischen Korrektur werden $N=\monopolN$ aufeinanderfolgende nichttriviale
Nullstellen $\{\gamma_n\}$ ab $t=\monopolTmin$ mit adaptiver \texttt{mpmath.siegelz}-Pipeline
bestimmt (Laufzeit des Referenzlaufs: $106{,}6\,\mathrm{s}$).
Die Entfaltung (Unfolding) folgt dem glatten Weyl-Hauptterm
\begin{equation}
  N(t) = \frac{t}{2\pi}\,\log\!\Bigl(\frac{t}{2\pi e}\Bigr) + \frac{7}{8},
  \label{eq:weyl-unfold}
\end{equation}
und liefert $N-1$ normierte Abstände $s_n = N(\gamma_{n+1}) - N(\gamma_n)$.
Als Referenz für das zweite Moment dient die GUE-Varianz
$\sigma^2_{\mathrm{GUE}} = \monopolGUE$ der Wigner-Dyson-Verteilung
(Poisson--GUE-Übergang, asymptotisch für große Matrizen).

\subsection{Ergebnisse bei $N=500$}

Im untersuchten Sektor $t \in [\monopolTmin,\,\monopolTmax]$ ergeben sich:

\begin{itemize}
  \item \textbf{Erstes Moment (Mittelwert):}
        $\bar{s} = \monopolMeanSpacing \approx 1$ --- konsistent mit Weyl-Konformität;
        der Monopol-Einfluss manifestiert sich hier \emph{nicht} als systematische
        Verschiebung des mittleren Abstands.
  \item \textbf{Zweites Moment (Varianz):}
        $\sigma^2 = \monopolVariance$;
        Abweichung vom GUE-Soll
        $\Delta\sigma^2 = \sigma^2 - \sigma^2_{\mathrm{GUE}} = \monopolDeltaSigma
        \approx +0{,}054$.
  \item \textbf{Standardabweichung:}
        $\monopolStdSpacing$.
\end{itemize}

Die beobachtete Varianz liegt damit etwa $51\,\%$ über dem GUE-Referenzwert
($\monopolVariance / \monopolGUE \approx 1{,}52$).
Dieser Überschuss im \emph{zweiten Moment} ist die numerisch gesicherte Größe,
an der der Monopol-Kopplungseffekt im Preprint verankert wird --- nicht eine
Verschiebung des mittleren Abstands um einen Faktor $\sim 3$.

\subsection{Monopol-Modellgrößen und spektraler Druck}

Im EABC-Monopolmodell bleibt das dominante Punkt-Monopol mit
$Q_{\mathrm{tot}} = \monopolQtot$ und dem daraus abgeleiteten Skalierungsfaktor
\begin{equation}
  f_{\mathrm{mon}} = \frac{Q_{\mathrm{tot}}}{\sqrt{3}} \approx \monopolFactor
  \label{eq:monopol-factor}
\end{equation}
als strukturelle Kopplungskonstante bestehen.
Als aggregierte Kopplungsgröße wird der \emph{spektrale Monopol-Druck}
\begin{equation}
  P_N = \sum_{n=1}^{N} \sin\!\Bigl(\gamma_n \,\frac{\log f_{\mathrm{mon}}}{150}\Bigr)
  \label{eq:monopol-pressure}
\end{equation}
aus dem numerischen Testskript berechnet.
Für $N=\monopolN$ ergibt sich $P_N = \monopolPressure$.

Da $P_N$ in \eqref{eq:monopol-pressure} als \emph{Summe} oszillierender Terme definiert ist,
skaliert $|P_N|$ nicht linear mit $N$; die Phase der einzelnen $\gamma_n$ bestimmt
konstruktive bzw.\ destruktive Interferenz
(z.\,B.\ $P_{50}\approx 2{,}56$, $P_{500}=\monopolPressure$).
Für vergleichende Angaben zwischen verschiedenen Stichproben ist daher stets die
\emph{Stichprobengröße $N$} mitzuführen.
Als heuristische Intensitätsgröße kann $P_N/N$ herangezogen werden
($P_{500}/500 \approx \monopolPressureNorm$); wegen der Oszillation ist $P_N/N$
für kleine $N$ nicht stabil und wird hier nicht als Konvergenzkennzahl verwendet.

\subsection{Konvergenz der Varianz}

Tabelle~\ref{tab:varianz-konvergenz} fasst Teilstichproben desselben Referenzlaufs
zusammen. Der mittlere Abstand bleibt durchgängig nahe $1$; $\Delta\sigma^2$ stabilisiert
sich im Bereich $+0{,}05$ bis $+0{,}07$ über $N$.

\begin{table}[htbp]
\centering
\caption{Konvergenz normierter Abstandsstatistiken (gleicher Lauf, $t_0=5000$).
$\sigma^2_{\mathrm{GUE}} = \monopolGUE$.}
\label{tab:varianz-konvergenz}
\begin{tabular}{rccc}
\toprule
$N$ & $\bar{s}$ & $\sigma^2$ & $\Delta\sigma^2$ \\
\midrule
50   & $\convFiftyMean$      & $\convFiftyVar$      & $+\convFiftyDelta$      \\
100  & $\convHundredMean$     & $\convHundredVar$     & $+\convHundredDelta$     \\
200  & $\convTwoHundredMean$  & $\convTwoHundredVar$  & $+\convTwoHundredDelta$  \\
500  & $\monopolMeanSpacing$  & $\monopolVariance$    & $+\monopolDeltaSigma$    \\
\bottomrule
\end{tabular}
\end{table}

\subsection{Interpretationsrahmen (ohne Spekulation)}

Zusammenfassend gilt nach Eliminierung des numerischen Artefakts:

\begin{enumerate}
  \item Die \emph{globale} Nullstellenzählung bleibt Weyl-konform
        ($\bar{s}\approx 1$); eine Dichte-Reduktion um den Faktor $1/3$ ist
        nicht belegt.
  \item Eine \emph{messbare} Abweichung vom GUE-Referenzensemble zeigt sich im
        zweiten Moment der Abstandsverteilung ($\Delta\sigma^2 > 0$ bei
        $N=\monopolN$).
  \item Der Monopol-Kopplungsterm \eqref{eq:monopol-pressure} liefert eine
        reelle, stichprobenabhängige Aggregatgröße; bei Vergleichen ist stets
        $N$ mit anzugeben, da $P_N$ als oszillierende Summe nicht
        trivial mit $N$ skaliert.
\end{enumerate}

Eine weitergehende Zuordnung von $\Delta\sigma^2$ zu konkreten Feldtheorie-Mechanismen
(``Stark-Effekt'', Symmetriebrechung der Dichte usw.) wird bewusst zurückgestellt,
bis unabhängige Reproduktionen und größere Stichproben vorliegen.

% Voraussetzungen: amsmath, booktabs (ggf. hyperref für \ref)

\input{monopol_preprint_results.txt}

% Modellkonstanten (nicht laufzeitabhängig)
\newcommand{\monopolQtot}{1796}
\newcommand{\monopolFactor}{1036.92} % 1796/\sqrt{3}
\newcommand{\monopolPressureNorm}{0.032292} % \monopolPressure / \monopolN

% Konvergenzdaten (Teilstichproben, gleicher Lauf wie oben)
\newcommand{\convFiftyVar}{0.172806}
\newcommand{\convFiftyDelta}{0.068616}
\newcommand{\convFiftyMean}{0.992296}
\newcommand{\convHundredVar}{0.155540}
\newcommand{\convHundredDelta}{0.051350}
\newcommand{\convHundredMean}{0.995738}
\newcommand{\convTwoHundredVar}{0.160372}
\newcommand{\convTwoHundredDelta}{0.056182}
\newcommand{\convTwoHundredMean}{0.998893}

% Diagnosewerte der fehlerhaften Hauptterm-Pipeline (Referenzlauf, $t\in[5000,5140]$)
\newcommand{\badMeanFixed}{2.9586}
\newcommand{\badVarFixed}{7.33}
\newcommand{\badMeanAdapt}{2.9962}
\newcommand{\badVarAdapt}{7.88}
\newcommand{\badCountInterval}{52}
\newcommand{\goodCountInterval}{150}

%=============================================================================
\section{Methodische Richtigstellung zur numerischen Nullstellen-Asymptotik bei hohen Parametern}
\label{sec:methodische-richtigstellung}
%=============================================================================

In der numerischen Verifikation des EABC-Monopolmodells im hochenergetischen Bereich
($t \approx 5000$) zeigt sich eine methodische Divergenz, die für jede spektrale
Auswertung zuerst geklärt werden muss.
Eine Auswertung der Riemann-Siegel-Z-Funktion, die nur den asymptotischen Hauptterm
verwendet und den Restterm $R(t)$ vernachlässigt, führt bei der Nullstellensuche
über Vorzeichenwechsel zu systematischen Lücken in der Stichprobe.

\begin{table}[htbp]
\centering
\caption{Vergleich numerischer Approximationsmethoden im Intervall $t \in [5000,\,5140]$.
Die ersten beiden Zeilen verwenden den unvollständigen Hauptterm; die dritte Zeile die
vollständige Implementierung \texttt{mpmath.siegelz} mit adaptiver Schrittweite
und \texttt{brentq}-Verfeinerung (Stichprobe $N=50$).}
\label{tab:siegel_vergleich}
\begin{tabular}{lccc}
\toprule
\textbf{Methode} &
\textbf{Mittlerer norm.\ Abstand} &
\textbf{Varianz $\sigma^2$} &
\textbf{Gefundene Nullstellen} \\
\midrule
Hauptterm, Schrittweite $0{,}5$ &
$\badMeanFixed$ & $\badVarFixed$ & $\sim\badCountInterval$ \\
Hauptterm, adaptive Schrittweite &
$\badMeanAdapt$ & $\badVarAdapt$ & $\sim\badCountInterval$ \\
\texttt{mpmath.siegelz}, adaptiv &
$\convFiftyMean$ & $\convFiftyVar$ & $\sim\goodCountInterval$ \\
\bottomrule
\end{tabular}
\end{table}

Tabelle~\ref{tab:siegel_vergleich} zeigt: Eine rein adaptive Schrittweite behebt das
Problem nicht, solange $R(t)$ fehlt---es werden dieselben $\sim\badCountInterval$ Nullstellen
gefunden wie mit fester Schrittweite $0{,}5$.
Der scheinbare Faktor $\sim 3$ im mittleren normierten Abstand ist damit ein
Aliasing-Effekt der Unterabtastung, kein spektraler Beitrag des Monopols.

Erst unter Einbeziehung der vollständigen Riemann-Siegel-Asymptotik stellt sich
Weyl-konformes Verhalten mit mittlerem Abstand $\approx 1$ ein.
Eine Modifikation der globalen Zählfunktion der Form
$N_{\mathrm{EABC}}(T) \approx \tfrac{1}{3}\,N_{\mathrm{standard}}(T)$
--- oder jede daraus abgeleitete ``Drittel-Symmetrie'' der Nullstellendichte ---
ist ohne unabhängige Evidenz nicht haltbar und wird verworfen.

%=============================================================================
\section{Spektrale Abweichungsanalyse und Monopol-Kopplung}
\label{sec:spektrale-abweichung}
%=============================================================================

\subsection{Statistisches Setup}

Nach der methodischen Korrektur werden $N=\monopolN$ aufeinanderfolgende nichttriviale
Nullstellen $\{\gamma_n\}$ ab $t=\monopolTmin$ mit adaptiver \texttt{mpmath.siegelz}-Pipeline
bestimmt (Laufzeit des Referenzlaufs: $106{,}6\,\mathrm{s}$).
Die Entfaltung (Unfolding) folgt dem glatten Weyl-Hauptterm
\begin{equation}
  N(t) = \frac{t}{2\pi}\,\log\!\Bigl(\frac{t}{2\pi e}\Bigr) + \frac{7}{8},
  \label{eq:weyl-unfold}
\end{equation}
und liefert $N-1$ normierte Abstände $s_n = N(\gamma_{n+1}) - N(\gamma_n)$.
Als Referenz für das zweite Moment dient die GUE-Varianz
$\sigma^2_{\mathrm{GUE}} = \monopolGUE$ der Wigner-Dyson-Verteilung
(Poisson--GUE-Übergang, asymptotisch für große Matrizen).

\subsection{Ergebnisse bei $N=500$}

Im untersuchten Sektor $t \in [\monopolTmin,\,\monopolTmax]$ ergeben sich:

\begin{itemize}
  \item \textbf{Erstes Moment (Mittelwert):}
        $\bar{s} = \monopolMeanSpacing \approx 1$ --- konsistent mit Weyl-Konformität;
        der Monopol-Einfluss manifestiert sich hier \emph{nicht} als systematische
        Verschiebung des mittleren Abstands.
  \item \textbf{Zweites Moment (Varianz):}
        $\sigma^2 = \monopolVariance$;
        Abweichung vom GUE-Soll
        $\Delta\sigma^2 = \sigma^2 - \sigma^2_{\mathrm{GUE}} = \monopolDeltaSigma
        \approx +0{,}054$.
  \item \textbf{Standardabweichung:}
        $\monopolStdSpacing$.
\end{itemize}

Die beobachtete Varianz liegt damit etwa $51\,\%$ über dem GUE-Referenzwert
($\monopolVariance / \monopolGUE \approx 1{,}52$).
Dieser Überschuss im \emph{zweiten Moment} ist die numerisch gesicherte Größe,
an der der Monopol-Kopplungseffekt im Preprint verankert wird --- nicht eine
Verschiebung des mittleren Abstands um einen Faktor $\sim 3$.

\subsection{Monopol-Modellgrößen und spektraler Druck}

Im EABC-Monopolmodell bleibt das dominante Punkt-Monopol mit
$Q_{\mathrm{tot}} = \monopolQtot$ und dem daraus abgeleiteten Skalierungsfaktor
\begin{equation}
  f_{\mathrm{mon}} = \frac{Q_{\mathrm{tot}}}{\sqrt{3}} \approx \monopolFactor
  \label{eq:monopol-factor}
\end{equation}
als strukturelle Kopplungskonstante bestehen.
Als aggregierte Kopplungsgröße wird der \emph{spektrale Monopol-Druck}
\begin{equation}
  P_N = \sum_{n=1}^{N} \sin\!\Bigl(\gamma_n \,\frac{\log f_{\mathrm{mon}}}{150}\Bigr)
  \label{eq:monopol-pressure}
\end{equation}
aus dem numerischen Testskript berechnet.
Für $N=\monopolN$ ergibt sich $P_N = \monopolPressure$.

Da $P_N$ in \eqref{eq:monopol-pressure} als \emph{Summe} oszillierender Terme definiert ist,
skaliert $|P_N|$ nicht linear mit $N$; die Phase der einzelnen $\gamma_n$ bestimmt
konstruktive bzw.\ destruktive Interferenz
(z.\,B.\ $P_{50}\approx 2{,}56$, $P_{500}=\monopolPressure$).
Für vergleichende Angaben zwischen verschiedenen Stichproben ist daher stets die
\emph{Stichprobengröße $N$} mitzuführen.
Als heuristische Intensitätsgröße kann $P_N/N$ herangezogen werden
($P_{500}/500 \approx \monopolPressureNorm$); wegen der Oszillation ist $P_N/N$
für kleine $N$ nicht stabil und wird hier nicht als Konvergenzkennzahl verwendet.

\subsection{Konvergenz der Varianz}

Tabelle~\ref{tab:varianz-konvergenz} fasst Teilstichproben desselben Referenzlaufs
zusammen. Der mittlere Abstand bleibt durchgängig nahe $1$; $\Delta\sigma^2$ stabilisiert
sich im Bereich $+0{,}05$ bis $+0{,}07$ über $N$.

\begin{table}[htbp]
\centering
\caption{Konvergenz normierter Abstandsstatistiken (gleicher Lauf, $t_0=5000$).
$\sigma^2_{\mathrm{GUE}} = \monopolGUE$.}
\label{tab:varianz-konvergenz}
\begin{tabular}{rccc}
\toprule
$N$ & $\bar{s}$ & $\sigma^2$ & $\Delta\sigma^2$ \\
\midrule
50   & $\convFiftyMean$      & $\convFiftyVar$      & $+\convFiftyDelta$      \\
100  & $\convHundredMean$     & $\convHundredVar$     & $+\convHundredDelta$     \\
200  & $\convTwoHundredMean$  & $\convTwoHundredVar$  & $+\convTwoHundredDelta$  \\
500  & $\monopolMeanSpacing$  & $\monopolVariance$    & $+\monopolDeltaSigma$    \\
\bottomrule
\end{tabular}
\end{table}

\subsection{Interpretationsrahmen (ohne Spekulation)}

Zusammenfassend gilt nach Eliminierung des numerischen Artefakts:

\begin{enumerate}
  \item Die \emph{globale} Nullstellenzählung bleibt Weyl-konform
        ($\bar{s}\approx 1$); eine Dichte-Reduktion um den Faktor $1/3$ ist
        nicht belegt.
  \item Eine \emph{messbare} Abweichung vom GUE-Referenzensemble zeigt sich im
        zweiten Moment der Abstandsverteilung ($\Delta\sigma^2 > 0$ bei
        $N=\monopolN$).
  \item Der Monopol-Kopplungsterm \eqref{eq:monopol-pressure} liefert eine
        reelle, stichprobenabhängige Aggregatgröße; bei Vergleichen ist stets
        $N$ mit anzugeben, da $P_N$ als oszillierende Summe nicht
        trivial mit $N$ skaliert.
\end{enumerate}

Eine weitergehende Zuordnung von $\Delta\sigma^2$ zu konkreten Feldtheorie-Mechanismen
(``Stark-Effekt'', Symmetriebrechung der Dichte usw.) wird bewusst zurückgestellt,
bis unabhängige Reproduktionen und größere Stichproben vorliegen.

% Voraussetzungen: amsmath, booktabs (ggf. hyperref für \ref)

\input{monopol_preprint_results.txt}

% Modellkonstanten (nicht laufzeitabhängig)
\newcommand{\monopolQtot}{1796}
\newcommand{\monopolFactor}{1036.92} % 1796/\sqrt{3}
\newcommand{\monopolPressureNorm}{0.032292} % \monopolPressure / \monopolN

% Konvergenzdaten (Teilstichproben, gleicher Lauf wie oben)
\newcommand{\convFiftyVar}{0.172806}
\newcommand{\convFiftyDelta}{0.068616}
\newcommand{\convFiftyMean}{0.992296}
\newcommand{\convHundredVar}{0.155540}
\newcommand{\convHundredDelta}{0.051350}
\newcommand{\convHundredMean}{0.995738}
\newcommand{\convTwoHundredVar}{0.160372}
\newcommand{\convTwoHundredDelta}{0.056182}
\newcommand{\convTwoHundredMean}{0.998893}

% Diagnosewerte der fehlerhaften Hauptterm-Pipeline (Referenzlauf, $t\in[5000,5140]$)
\newcommand{\badMeanFixed}{2.9586}
\newcommand{\badVarFixed}{7.33}
\newcommand{\badMeanAdapt}{2.9962}
\newcommand{\badVarAdapt}{7.88}
\newcommand{\badCountInterval}{52}
\newcommand{\goodCountInterval}{150}

%=============================================================================
\section{Methodische Richtigstellung zur numerischen Nullstellen-Asymptotik bei hohen Parametern}
\label{sec:methodische-richtigstellung}
%=============================================================================

In der numerischen Verifikation des EABC-Monopolmodells im hochenergetischen Bereich
($t \approx 5000$) zeigt sich eine methodische Divergenz, die für jede spektrale
Auswertung zuerst geklärt werden muss.
Eine Auswertung der Riemann-Siegel-Z-Funktion, die nur den asymptotischen Hauptterm
verwendet und den Restterm $R(t)$ vernachlässigt, führt bei der Nullstellensuche
über Vorzeichenwechsel zu systematischen Lücken in der Stichprobe.

\begin{table}[htbp]
\centering
\caption{Vergleich numerischer Approximationsmethoden im Intervall $t \in [5000,\,5140]$.
Die ersten beiden Zeilen verwenden den unvollständigen Hauptterm; die dritte Zeile die
vollständige Implementierung \texttt{mpmath.siegelz} mit adaptiver Schrittweite
und \texttt{brentq}-Verfeinerung (Stichprobe $N=50$).}
\label{tab:siegel_vergleich}
\begin{tabular}{lccc}
\toprule
\textbf{Methode} &
\textbf{Mittlerer norm.\ Abstand} &
\textbf{Varianz $\sigma^2$} &
\textbf{Gefundene Nullstellen} \\
\midrule
Hauptterm, Schrittweite $0{,}5$ &
$\badMeanFixed$ & $\badVarFixed$ & $\sim\badCountInterval$ \\
Hauptterm, adaptive Schrittweite &
$\badMeanAdapt$ & $\badVarAdapt$ & $\sim\badCountInterval$ \\
\texttt{mpmath.siegelz}, adaptiv &
$\convFiftyMean$ & $\convFiftyVar$ & $\sim\goodCountInterval$ \\
\bottomrule
\end{tabular}
\end{table}

Tabelle~\ref{tab:siegel_vergleich} zeigt: Eine rein adaptive Schrittweite behebt das
Problem nicht, solange $R(t)$ fehlt---es werden dieselben $\sim\badCountInterval$ Nullstellen
gefunden wie mit fester Schrittweite $0{,}5$.
Der scheinbare Faktor $\sim 3$ im mittleren normierten Abstand ist damit ein
Aliasing-Effekt der Unterabtastung, kein spektraler Beitrag des Monopols.

Erst unter Einbeziehung der vollständigen Riemann-Siegel-Asymptotik stellt sich
Weyl-konformes Verhalten mit mittlerem Abstand $\approx 1$ ein.
Eine Modifikation der globalen Zählfunktion der Form
$N_{\mathrm{EABC}}(T) \approx \tfrac{1}{3}\,N_{\mathrm{standard}}(T)$
--- oder jede daraus abgeleitete ``Drittel-Symmetrie'' der Nullstellendichte ---
ist ohne unabhängige Evidenz nicht haltbar und wird verworfen.

%=============================================================================
\section{Spektrale Abweichungsanalyse und Monopol-Kopplung}
\label{sec:spektrale-abweichung}
%=============================================================================

\subsection{Statistisches Setup}

Nach der methodischen Korrektur werden $N=\monopolN$ aufeinanderfolgende nichttriviale
Nullstellen $\{\gamma_n\}$ ab $t=\monopolTmin$ mit adaptiver \texttt{mpmath.siegelz}-Pipeline
bestimmt (Laufzeit des Referenzlaufs: $106{,}6\,\mathrm{s}$).
Die Entfaltung (Unfolding) folgt dem glatten Weyl-Hauptterm
\begin{equation}
  N(t) = \frac{t}{2\pi}\,\log\!\Bigl(\frac{t}{2\pi e}\Bigr) + \frac{7}{8},
  \label{eq:weyl-unfold}
\end{equation}
und liefert $N-1$ normierte Abstände $s_n = N(\gamma_{n+1}) - N(\gamma_n)$.
Als Referenz für das zweite Moment dient die GUE-Varianz
$\sigma^2_{\mathrm{GUE}} = \monopolGUE$ der Wigner-Dyson-Verteilung
(Poisson--GUE-Übergang, asymptotisch für große Matrizen).

\subsection{Ergebnisse bei $N=500$}

Im untersuchten Sektor $t \in [\monopolTmin,\,\monopolTmax]$ ergeben sich:

\begin{itemize}
  \item \textbf{Erstes Moment (Mittelwert):}
        $\bar{s} = \monopolMeanSpacing \approx 1$ --- konsistent mit Weyl-Konformität;
        der Monopol-Einfluss manifestiert sich hier \emph{nicht} als systematische
        Verschiebung des mittleren Abstands.
  \item \textbf{Zweites Moment (Varianz):}
        $\sigma^2 = \monopolVariance$;
        Abweichung vom GUE-Soll
        $\Delta\sigma^2 = \sigma^2 - \sigma^2_{\mathrm{GUE}} = \monopolDeltaSigma
        \approx +0{,}054$.
  \item \textbf{Standardabweichung:}
        $\monopolStdSpacing$.
\end{itemize}

Die beobachtete Varianz liegt damit etwa $51\,\%$ über dem GUE-Referenzwert
($\monopolVariance / \monopolGUE \approx 1{,}52$).
Dieser Überschuss im \emph{zweiten Moment} ist die numerisch gesicherte Größe,
an der der Monopol-Kopplungseffekt im Preprint verankert wird --- nicht eine
Verschiebung des mittleren Abstands um einen Faktor $\sim 3$.

\subsection{Monopol-Modellgrößen und spektraler Druck}

Im EABC-Monopolmodell bleibt das dominante Punkt-Monopol mit
$Q_{\mathrm{tot}} = \monopolQtot$ und dem daraus abgeleiteten Skalierungsfaktor
\begin{equation}
  f_{\mathrm{mon}} = \frac{Q_{\mathrm{tot}}}{\sqrt{3}} \approx \monopolFactor
  \label{eq:monopol-factor}
\end{equation}
als strukturelle Kopplungskonstante bestehen.
Als aggregierte Kopplungsgröße wird der \emph{spektrale Monopol-Druck}
\begin{equation}
  P_N = \sum_{n=1}^{N} \sin\!\Bigl(\gamma_n \,\frac{\log f_{\mathrm{mon}}}{150}\Bigr)
  \label{eq:monopol-pressure}
\end{equation}
aus dem numerischen Testskript berechnet.
Für $N=\monopolN$ ergibt sich $P_N = \monopolPressure$.

Da $P_N$ in \eqref{eq:monopol-pressure} als \emph{Summe} oszillierender Terme definiert ist,
skaliert $|P_N|$ nicht linear mit $N$; die Phase der einzelnen $\gamma_n$ bestimmt
konstruktive bzw.\ destruktive Interferenz
(z.\,B.\ $P_{50}\approx 2{,}56$, $P_{500}=\monopolPressure$).
Für vergleichende Angaben zwischen verschiedenen Stichproben ist daher stets die
\emph{Stichprobengröße $N$} mitzuführen.
Als heuristische Intensitätsgröße kann $P_N/N$ herangezogen werden
($P_{500}/500 \approx \monopolPressureNorm$); wegen der Oszillation ist $P_N/N$
für kleine $N$ nicht stabil und wird hier nicht als Konvergenzkennzahl verwendet.

\subsection{Konvergenz der Varianz}

Tabelle~\ref{tab:varianz-konvergenz} fasst Teilstichproben desselben Referenzlaufs
zusammen. Der mittlere Abstand bleibt durchgängig nahe $1$; $\Delta\sigma^2$ stabilisiert
sich im Bereich $+0{,}05$ bis $+0{,}07$ über $N$.

\begin{table}[htbp]
\centering
\caption{Konvergenz normierter Abstandsstatistiken (gleicher Lauf, $t_0=5000$).
$\sigma^2_{\mathrm{GUE}} = \monopolGUE$.}
\label{tab:varianz-konvergenz}
\begin{tabular}{rccc}
\toprule
$N$ & $\bar{s}$ & $\sigma^2$ & $\Delta\sigma^2$ \\
\midrule
50   & $\convFiftyMean$      & $\convFiftyVar$      & $+\convFiftyDelta$      \\
100  & $\convHundredMean$     & $\convHundredVar$     & $+\convHundredDelta$     \\
200  & $\convTwoHundredMean$  & $\convTwoHundredVar$  & $+\convTwoHundredDelta$  \\
500  & $\monopolMeanSpacing$  & $\monopolVariance$    & $+\monopolDeltaSigma$    \\
\bottomrule
\end{tabular}
\end{table}

\subsection{Interpretationsrahmen (ohne Spekulation)}

Zusammenfassend gilt nach Eliminierung des numerischen Artefakts:

\begin{enumerate}
  \item Die \emph{globale} Nullstellenzählung bleibt Weyl-konform
        ($\bar{s}\approx 1$); eine Dichte-Reduktion um den Faktor $1/3$ ist
        nicht belegt.
  \item Eine \emph{messbare} Abweichung vom GUE-Referenzensemble zeigt sich im
        zweiten Moment der Abstandsverteilung ($\Delta\sigma^2 > 0$ bei
        $N=\monopolN$).
  \item Der Monopol-Kopplungsterm \eqref{eq:monopol-pressure} liefert eine
        reelle, stichprobenabhängige Aggregatgröße; bei Vergleichen ist stets
        $N$ mit anzugeben, da $P_N$ als oszillierende Summe nicht
        trivial mit $N$ skaliert.
\end{enumerate}

Eine weitergehende Zuordnung von $\Delta\sigma^2$ zu konkreten Feldtheorie-Mechanismen
(``Stark-Effekt'', Symmetriebrechung der Dichte usw.) wird bewusst zurückgestellt,
bis unabhängige Reproduktionen und größere Stichproben vorliegen.

% Voraussetzungen: amsmath, booktabs (ggf. hyperref für \ref)

\input{monopol_preprint_results.txt}

% Modellkonstanten (nicht laufzeitabhängig)
\newcommand{\monopolQtot}{1796}
\newcommand{\monopolFactor}{1036.92} % 1796/\sqrt{3}
\newcommand{\monopolPressureNorm}{0.032292} % \monopolPressure / \monopolN

% Konvergenzdaten (Teilstichproben, gleicher Lauf wie oben)
\newcommand{\convFiftyVar}{0.172806}
\newcommand{\convFiftyDelta}{0.068616}
\newcommand{\convFiftyMean}{0.992296}
\newcommand{\convHundredVar}{0.155540}
\newcommand{\convHundredDelta}{0.051350}
\newcommand{\convHundredMean}{0.995738}
\newcommand{\convTwoHundredVar}{0.160372}
\newcommand{\convTwoHundredDelta}{0.056182}
\newcommand{\convTwoHundredMean}{0.998893}

% Diagnosewerte der fehlerhaften Hauptterm-Pipeline (Referenzlauf, $t\in[5000,5140]$)
\newcommand{\badMeanFixed}{2.9586}
\newcommand{\badVarFixed}{7.33}
\newcommand{\badMeanAdapt}{2.9962}
\newcommand{\badVarAdapt}{7.88}
\newcommand{\badCountInterval}{52}
\newcommand{\goodCountInterval}{150}

%=============================================================================
\section{Methodische Richtigstellung zur numerischen Nullstellen-Asymptotik bei hohen Parametern}
\label{sec:methodische-richtigstellung}
%=============================================================================

In der numerischen Verifikation des EABC-Monopolmodells im hochenergetischen Bereich
($t \approx 5000$) zeigt sich eine methodische Divergenz, die für jede spektrale
Auswertung zuerst geklärt werden muss.
Eine Auswertung der Riemann-Siegel-Z-Funktion, die nur den asymptotischen Hauptterm
verwendet und den Restterm $R(t)$ vernachlässigt, führt bei der Nullstellensuche
über Vorzeichenwechsel zu systematischen Lücken in der Stichprobe.

\begin{table}[htbp]
\centering
\caption{Vergleich numerischer Approximationsmethoden im Intervall $t \in [5000,\,5140]$.
Die ersten beiden Zeilen verwenden den unvollständigen Hauptterm; die dritte Zeile die
vollständige Implementierung \texttt{mpmath.siegelz} mit adaptiver Schrittweite
und \texttt{brentq}-Verfeinerung (Stichprobe $N=50$).}
\label{tab:siegel_vergleich}
\begin{tabular}{lccc}
\toprule
\textbf{Methode} &
\textbf{Mittlerer norm.\ Abstand} &
\textbf{Varianz $\sigma^2$} &
\textbf{Gefundene Nullstellen} \\
\midrule
Hauptterm, Schrittweite $0{,}5$ &
$\badMeanFixed$ & $\badVarFixed$ & $\sim\badCountInterval$ \\
Hauptterm, adaptive Schrittweite &
$\badMeanAdapt$ & $\badVarAdapt$ & $\sim\badCountInterval$ \\
\texttt{mpmath.siegelz}, adaptiv &
$\convFiftyMean$ & $\convFiftyVar$ & $\sim\goodCountInterval$ \\
\bottomrule
\end{tabular}
\end{table}

Tabelle~\ref{tab:siegel_vergleich} zeigt: Eine rein adaptive Schrittweite behebt das
Problem nicht, solange $R(t)$ fehlt---es werden dieselben $\sim\badCountInterval$ Nullstellen
gefunden wie mit fester Schrittweite $0{,}5$.
Der scheinbare Faktor $\sim 3$ im mittleren normierten Abstand ist damit ein
Aliasing-Effekt der Unterabtastung, kein spektraler Beitrag des Monopols.

Erst unter Einbeziehung der vollständigen Riemann-Siegel-Asymptotik stellt sich
Weyl-konformes Verhalten mit mittlerem Abstand $\approx 1$ ein.
Eine Modifikation der globalen Zählfunktion der Form
$N_{\mathrm{EABC}}(T) \approx \tfrac{1}{3}\,N_{\mathrm{standard}}(T)$
--- oder jede daraus abgeleitete ``Drittel-Symmetrie'' der Nullstellendichte ---
ist ohne unabhängige Evidenz nicht haltbar und wird verworfen.

%=============================================================================
\section{Spektrale Abweichungsanalyse und Monopol-Kopplung}
\label{sec:spektrale-abweichung}
%=============================================================================

\subsection{Statistisches Setup}

Nach der methodischen Korrektur werden $N=\monopolN$ aufeinanderfolgende nichttriviale
Nullstellen $\{\gamma_n\}$ ab $t=\monopolTmin$ mit adaptiver \texttt{mpmath.siegelz}-Pipeline
bestimmt (Laufzeit des Referenzlaufs: $106{,}6\,\mathrm{s}$).
Die Entfaltung (Unfolding) folgt dem glatten Weyl-Hauptterm
\begin{equation}
  N(t) = \frac{t}{2\pi}\,\log\!\Bigl(\frac{t}{2\pi e}\Bigr) + \frac{7}{8},
  \label{eq:weyl-unfold}
\end{equation}
und liefert $N-1$ normierte Abstände $s_n = N(\gamma_{n+1}) - N(\gamma_n)$.
Als Referenz für das zweite Moment dient die GUE-Varianz
$\sigma^2_{\mathrm{GUE}} = \monopolGUE$ der Wigner-Dyson-Verteilung
(Poisson--GUE-Übergang, asymptotisch für große Matrizen).

\subsection{Ergebnisse bei $N=500$}

Im untersuchten Sektor $t \in [\monopolTmin,\,\monopolTmax]$ ergeben sich:

\begin{itemize}
  \item \textbf{Erstes Moment (Mittelwert):}
        $\bar{s} = \monopolMeanSpacing \approx 1$ --- konsistent mit Weyl-Konformität;
        der Monopol-Einfluss manifestiert sich hier \emph{nicht} als systematische
        Verschiebung des mittleren Abstands.
  \item \textbf{Zweites Moment (Varianz):}
        $\sigma^2 = \monopolVariance$;
        Abweichung vom GUE-Soll
        $\Delta\sigma^2 = \sigma^2 - \sigma^2_{\mathrm{GUE}} = \monopolDeltaSigma
        \approx +0{,}054$.
  \item \textbf{Standardabweichung:}
        $\monopolStdSpacing$.
\end{itemize}

Die beobachtete Varianz liegt damit etwa $51\,\%$ über dem GUE-Referenzwert
($\monopolVariance / \monopolGUE \approx 1{,}52$).
Dieser Überschuss im \emph{zweiten Moment} ist die numerisch gesicherte Größe,
an der der Monopol-Kopplungseffekt im Preprint verankert wird --- nicht eine
Verschiebung des mittleren Abstands um einen Faktor $\sim 3$.

\subsection{Monopol-Modellgrößen und spektraler Druck}

Im EABC-Monopolmodell bleibt das dominante Punkt-Monopol mit
$Q_{\mathrm{tot}} = \monopolQtot$ und dem daraus abgeleiteten Skalierungsfaktor
\begin{equation}
  f_{\mathrm{mon}} = \frac{Q_{\mathrm{tot}}}{\sqrt{3}} \approx \monopolFactor
  \label{eq:monopol-factor}
\end{equation}
als strukturelle Kopplungskonstante bestehen.
Als aggregierte Kopplungsgröße wird der \emph{spektrale Monopol-Druck}
\begin{equation}
  P_N = \sum_{n=1}^{N} \sin\!\Bigl(\gamma_n \,\frac{\log f_{\mathrm{mon}}}{150}\Bigr)
  \label{eq:monopol-pressure}
\end{equation}
aus dem numerischen Testskript berechnet.
Für $N=\monopolN$ ergibt sich $P_N = \monopolPressure$.

Da $P_N$ in \eqref{eq:monopol-pressure} als \emph{Summe} oszillierender Terme definiert ist,
skaliert $|P_N|$ nicht linear mit $N$; die Phase der einzelnen $\gamma_n$ bestimmt
konstruktive bzw.\ destruktive Interferenz
(z.\,B.\ $P_{50}\approx 2{,}56$, $P_{500}=\monopolPressure$).
Für vergleichende Angaben zwischen verschiedenen Stichproben ist daher stets die
\emph{Stichprobengröße $N$} mitzuführen.
Als heuristische Intensitätsgröße kann $P_N/N$ herangezogen werden
($P_{500}/500 \approx \monopolPressureNorm$); wegen der Oszillation ist $P_N/N$
für kleine $N$ nicht stabil und wird hier nicht als Konvergenzkennzahl verwendet.

\subsection{Konvergenz der Varianz}

Tabelle~\ref{tab:varianz-konvergenz} fasst Teilstichproben desselben Referenzlaufs
zusammen. Der mittlere Abstand bleibt durchgängig nahe $1$; $\Delta\sigma^2$ stabilisiert
sich im Bereich $+0{,}05$ bis $+0{,}07$ über $N$.

\begin{table}[htbp]
\centering
\caption{Konvergenz normierter Abstandsstatistiken (gleicher Lauf, $t_0=5000$).
$\sigma^2_{\mathrm{GUE}} = \monopolGUE$.}
\label{tab:varianz-konvergenz}
\begin{tabular}{rccc}
\toprule
$N$ & $\bar{s}$ & $\sigma^2$ & $\Delta\sigma^2$ \\
\midrule
50   & $\convFiftyMean$      & $\convFiftyVar$      & $+\convFiftyDelta$      \\
100  & $\convHundredMean$     & $\convHundredVar$     & $+\convHundredDelta$     \\
200  & $\convTwoHundredMean$  & $\convTwoHundredVar$  & $+\convTwoHundredDelta$  \\
500  & $\monopolMeanSpacing$  & $\monopolVariance$    & $+\monopolDeltaSigma$    \\
\bottomrule
\end{tabular}
\end{table}

\subsection{Interpretationsrahmen (ohne Spekulation)}

Zusammenfassend gilt nach Eliminierung des numerischen Artefakts:

\begin{enumerate}
  \item Die \emph{globale} Nullstellenzählung bleibt Weyl-konform
        ($\bar{s}\approx 1$); eine Dichte-Reduktion um den Faktor $1/3$ ist
        nicht belegt.
  \item Eine \emph{messbare} Abweichung vom GUE-Referenzensemble zeigt sich im
        zweiten Moment der Abstandsverteilung ($\Delta\sigma^2 > 0$ bei
        $N=\monopolN$).
  \item Der Monopol-Kopplungsterm \eqref{eq:monopol-pressure} liefert eine
        reelle, stichprobenabhängige Aggregatgröße; bei Vergleichen ist stets
        $N$ mit anzugeben, da $P_N$ als oszillierende Summe nicht
        trivial mit $N$ skaliert.
\end{enumerate}

Eine weitergehende Zuordnung von $\Delta\sigma^2$ zu konkreten Feldtheorie-Mechanismen
(``Stark-Effekt'', Symmetriebrechung der Dichte usw.) wird bewusst zurückgestellt,
bis unabhängige Reproduktionen und größere Stichproben vorliegen.
